\documentclass[UTF8,12pt,a4paper]{ctexart}

\usepackage{amsmath,amssymb,bm}
\usepackage{booktabs}
\usepackage{geometry}
\usepackage{longtable}
\usepackage{array}
\usepackage{hyperref}
\usepackage{graphicx}
\usepackage{tikz}
\usetikzlibrary{arrows.meta,positioning}

\geometry{left=2.6cm,right=2.6cm,top=2.5cm,bottom=2.5cm}
\hypersetup{colorlinks=true,linkcolor=blue,citecolor=blue,urlcolor=blue}
\newcommand{\Rcal}{\mathcal{R}}
\newcommand{\Fsen}{\mathcal{F}}
\title{基于低速标定等效间隙梯度约束和投影局部搜索的叶尖间隙--振动联合识别方法}
\author{}
\date{}

\begin{document}

\maketitle

\begin{abstract}
径向间隙变化不仅会影响传统叶尖定时中的到达时刻提取，还会造成低速基准模板与高速完整波形失配，使间隙响应与周向振动在识别过程中相互混淆。针对这一问题，本文提出一种基于低速标定等效局部间隙梯度约束的叶尖间隙与振动联合识别方法。该方法首先由低速旋转数据建立传感器相关的实测基准模板，并结合静态多间隙响应面标定模板对应的局部间隙路径；随后在高速识别中固定该标定关系，以实测低速模板表示基准波形，以静态响应面描述相对于低速状态的径向间隙变化，并在完整波形模型中联合估计径向间隙增量和周向振动参数。该方法将高速识别前能够标定的局部间隙关系与需要由高速数据估计的参数分开，限制局部间隙梯度与振动参数在同一高速波形拟合中相互补偿。模型最终输出相对于低速标定状态的等效径向间隙变化，以及叶尖周向振动的频率、幅值和相位。本文给出多频一般形式，试验分析采用单主导频率特例，多频识别仍需通过仿真验证。
\end{abstract}

\noindent\textbf{关键词：}叶尖间隙；叶片振动；完整波形；等效间隙梯度；变量投影；联合识别

\section{引言}

叶尖间隙是反映旋转机械气动性能和运行安全裕度的重要参数。非接触式叶尖传感器能够记录叶尖经过探头敏感区域时的高速电压波形\cite{Yu2020,Muller1997}。与主要利用叶尖到达时刻或其相位信息的方法\cite{Diamond2021}相比，完整波形还包含幅值和局部形状信息，因此可为叶尖间隙和叶片振动的联合估计提供观测基础\cite{Heller2020}。

传统叶尖定时通常利用阈值交点、峰值位置或其他脉冲特征确定叶片到达时刻，并将其相对于无振动参考时刻的偏差换算为叶尖周向振动位移。径向间隙变化会改变传感器脉冲的幅值、宽度和局部形状，即使叶片真实周向位置不变，也可能使提取的到达时刻发生偏移，并将该偏差传递到后续振动参数识别中。

完整波形方法通过低速基准脉冲与高速实测脉冲之间的位置差异识别振动。高速径向间隙发生变化时，高速脉冲在发生周向位置变化的同时还会产生幅值和形状变化，因而不再是低速模板的单纯平移。若模型只允许低速模板发生位置移动，间隙变化产生的部分波形差异就可能被解释为周向振动位移，从而影响振动幅值和相位的识别。

为此，本文利用静态多间隙数据描述径向间隙变化引起的脉冲幅值和形状变化，利用低速旋转数据保存当前叶片与传感器组合的实测基准波形。低速模板与静态响应面共同确定局部间隙路径；高速阶段保持该局部关系不变，只估计相对于低速状态的径向间隙增量，并与周向振动参数联合求解。由此避免仅通过移动低速模板解释高速间隙变化产生的波形差异。

本文首先从叶尖定时和完整波形匹配两方面说明径向间隙变化对振动识别的影响，并给出基本建模假设。随后说明低速旋转数据和静态多间隙数据的功能分工，建立径向间隙增量与周向振动位移共同作用的完整波形模型，最后给出候选生成、投影局部搜索和完整非线性优化流程。

\section{研究问题与建模假设}

\subsection{叶尖间隙--振动联合识别问题}

本文研究的问题是：当高速运行状态的径向间隙不同于低速基准状态时，如何利用同一组叶尖传感器完整波形，同时识别相对径向间隙变化和叶尖周向振动，并避免将间隙变化引起的波形差异误认为振动位移。为明确该问题，设 $g$ 为叶尖与传感器之间沿敏感方向的等效间隙，在本文试验布置下近似对应径向间隙；设 $x$ 为叶尖经过传感器敏感区域时的局部周向坐标，$u(t)$ 为叶尖振动在周向上的等效位移。

传统叶尖定时根据阈值交点、峰值或相关定位等脉冲特征确定叶片到达时刻。设叶片周向振动对应的真实到达时刻偏差为 $\Delta t_u$，径向间隙变化引起的特征定位偏差为 $\delta t_g$，其他测量误差为 $\varepsilon_t$，则提取的到达时刻偏差可表示为
\begin{equation}
\widehat{\Delta t}=\Delta t_u+\delta t_g+\varepsilon_t.
\label{eq:btt_time_error}
\end{equation}
在局部周向速度 $V_{\mathrm{tip}}$ 已知的条件下，传统叶尖定时给出的周向位移为
\begin{equation}
\widehat u_{\mathrm{BTT}}
=V_{\mathrm{tip}}\widehat{\Delta t}
\approx u+V_{\mathrm{tip}}\delta t_g+V_{\mathrm{tip}}\varepsilon_t,
\label{eq:btt_displacement}
\end{equation}
其中，$u\approx V_{\mathrm{tip}}\Delta t_u$ 为真实周向振动位移。式\eqref{eq:btt_displacement}表明，径向间隙变化一旦使脉冲幅值、宽度或局部轮廓发生变化，并进一步改变特征定位结果，其影响就会以附加周向位移的形式进入振动识别。

完整波形方法不再依赖单一特征点，而是通过低速基准模板与高速实测脉冲的整体匹配估计周向位移，但间隙与振动的混淆并不会因此自动消失。为说明这一点，在单传感器、单脉冲和小扰动条件下，记低速基准模板为 $T(x)$，模板的空间导数为 $T_x(x)$。若无振动波形关于间隙和周向位置的局部响应记为 $v(g,x)$，则低速基准间隙路径 $g^{\mathrm{low}}(x)$ 附近的间隙灵敏度定义为
\begin{equation}
S_g(x)=
\left.
\frac{\partial v(g,x)}{\partial g}
\right|_{g=g^{\mathrm{low}}(x)}.
\label{eq:problem_gap_sensitivity}
\end{equation}
在当前脉冲窗口内，将周向位移 $u$ 和径向间隙增量 $\Delta g$ 视为给定的小量，并忽略二者的乘积及其他二阶项，则高速波形相对于低速模板的一阶变化可写为
\begin{equation}
v^{\mathrm{high}}(x)-T(x)
\approx -T_x(x)u+S_g(x)\Delta g.
\label{eq:problem_linearization}
\end{equation}
式\eqref{eq:problem_linearization}中，第一项表示周向位移引起的波形平移，第二项表示径向间隙变化引起的幅值和形状变化。若识别模型只允许固定模板发生周向平移，即用 $-T_x(x)\widetilde u$ 解释全部波形变化，则在当前窗口内有近似关系
\begin{equation}
\widetilde u-u
\approx
-\Delta g\,
\frac{\langle T_x,S_g\rangle}
{\langle T_x,T_x\rangle},
\label{eq:gap_vibration_bias}
\end{equation}
其中，$\langle\cdot,\cdot\rangle$ 表示在有效脉冲窗口内的离散内积，$\widetilde u$ 为固定模板模型估计的周向位移。只要间隙灵敏度 $S_g$ 在模板平移方向 $T_x$ 上具有非零投影，固定模板匹配就会把部分间隙响应分配给周向位移。本文需要解决的正是这种由波形灵敏度重叠造成的间隙--振动混淆，而不是单纯追求更小的电压拟合残差。式\eqref{eq:problem_linearization}和式\eqref{eq:gap_vibration_bias}仅用于说明研究问题，后续识别采用完整非线性波形模型。

因此，在静态多间隙标定数据、低速基准波形和高速完整波形均可获得的条件下，本文估计各有效传感器相对于低速状态的等效径向间隙增量，以及由多传感器共同约束的周向振动频率、幅值和参考相位。该问题属于低速标定约束下的局部联合识别，所得间隙量不解释为缺少独立间隙基准条件下的绝对间隙测量值。

\subsection{基本建模假设}

为使低速模板、静态间隙响应面和高速观测数据能够在同一前向模型中建立联系，本文采用以下基本假设：
\begin{enumerate}
  \item 低速工况下，目标频带内的叶片振动相对于目标识别精度可忽略，或其在多转聚合后形成的系统性分量可忽略；同一目标叶片的重复脉冲经统一坐标转换后具有良好重复性。因此，多转聚合波形可作为当前叶片--传感器组合的低速基准模板。
  \item 在所选分析窗口内，各传感器对应测点处的平均径向间隙变化时间尺度远大于目标振动周期。因此，每个测点相对于低速状态的径向变化可用窗口内恒定的间隙增量 $\Delta g_s$ 表示。
  \item 在本文研究工况和目标频带内，叶片振动对测量波形的主要动态作用表示为叶尖周向等效位移，高速与低速之间沿径向的主要差异表示为准静态间隙增量。由此，动态位移采用周向分量描述，径向方向用于描述工况之间的准静态间隙差异。
  \item 在单个传感器的有效敏感窗口内，叶尖局部几何与传感器安装关系从低速到高速保持足够稳定。因此，由低速模板与静态响应面联合标定得到的等效局部间隙路径可用于描述高速状态下周向位置变化所对应的局部等效间隙变化。
  \item 经信号极性、基线、电压尺度和空间坐标校正后，静态标定、低速旋转和高速旋转条件下的局部间隙--电压响应具有可比性。因此，静态响应面可用于计算高速状态相对于低速状态的间隙增量响应。
\end{enumerate}

\section{联合识别方法}

\subsection{方法总体框架}

本文方法按照数据形成和参数传递的先后顺序分为四个阶段，如图\ref{fig:method_framework}所示。第一阶段将低速和高速采样信号由时间域转换到统一的局部周向空间坐标，并将低速重复波形重采样到公共空间网格。第二阶段分别利用静态多间隙数据建立径向间隙--周向位置响应面，利用重采样后的低速旋转数据建立传感器相关的实测基准模板。第三阶段将低速模板与静态响应面进行有界匹配，确定低速等效基准间隙、等效局部间隙梯度以及空间和电压配准参数。第四阶段固定低速标定关系，以实测模板为基准、以静态响应面计算间隙增量响应，在完整波形模型中联合估计高速径向间隙增量和周向振动参数。

静态间隙库描述径向间隙变化引起的波形变化，低速模板保存当前叶片与传感器组合的实际基准波形。二者通过低速等效标定建立联系：响应面不替换实测模板，在高速阶段仅用于计算相对于低速状态的间隙增量响应。

\begin{figure}[htbp]
\centering
\resizebox{\textwidth}{!}{%
\begin{tikzpicture}[
  >=Latex,
  data/.style={draw,rounded corners,fill=gray!8,text width=33mm,minimum height=10mm,align=center},
  proc/.style={draw,rounded corners,text width=37mm,minimum height=11mm,align=center},
  resultbox/.style={draw,rounded corners,fill=gray!8,text width=39mm,minimum height=11mm,align=center},
  flow/.style={->,thick}
]
\node[data] (static) at (0,0) {静态多间隙数据\\$\{g_i,v_i^{\mathrm{stat}}\}$};
\node[data] (low) at (5.1,0) {低速旋转数据\\$v_{s\ell}^{\mathrm{low}}(t)$};
\node[data] (high) at (10.2,0) {高速旋转数据\\$v_s^{\mathrm{high}}(t)$};

\node[proc] (surface) at (0,-2.3) {空间坐标统一\\构建 $\Rcal(g,x)$};
\node[proc] (template) at (5.1,-2.3) {时间--空间转换与公共网格重采样\\构建 $T_s^{\mathrm{low}}(x)$};
\node[proc] (hsample) at (10.2,-2.3) {时间--空间转换\\形成 $\{t_{sn},x_{sn},v_{sn}\}$};

\node[proc] (cal) at (3.4,-4.7) {低速等效标定\\确定局部间隙路径与配准参数};
\node[proc] (joint) at (9.0,-4.7) {固定低速标定关系\\完整波形联合识别};
\node[resultbox] (result) at (9.0,-6.7) {相对径向间隙变化\\振动频率、幅值和相位};

\draw[flow] (static) -- (surface);
\draw[flow] (low) -- (template);
\draw[flow] (high) -- (hsample);
\draw[flow] (surface) |- (cal.west);
\draw[flow] (template) -- (cal.north);
\draw[flow] (cal) -- (joint);
\draw[flow] (hsample) -- (joint.north);
\draw[flow] (joint) -- (result);
\end{tikzpicture}%
}
\caption{叶尖间隙与振动联合识别方法总体流程}
\label{fig:method_framework}
\end{figure}

\subsection{时间--空间坐标转换与低速模板构建}

叶尖传感器首先在时间域记录电压序列，而静态响应面和完整波形模型均以叶尖相对于传感器的局部周向位置为自变量。因此，在构建低速模板和高速观测模型之前，需要利用转角参考将各转次的采样时刻转换为统一空间坐标。设相邻两个一次转速参考（once-per-revolution，OPR）时刻为 $t_m^{\mathrm{OPR}}$ 和 $t_{m+1}^{\mathrm{OPR}}$，两参考位置之间的已知转角增量为 $\Delta\theta_m^{\mathrm{OPR}}$。在该时间区间内，累计转角定义为
\begin{equation}
\theta(t)=\theta_m^{\mathrm{OPR}}
+\Delta\theta_m^{\mathrm{OPR}}
\frac{t-t_m^{\mathrm{OPR}}}
{t_{m+1}^{\mathrm{OPR}}-t_m^{\mathrm{OPR}}},
\qquad
t\in[t_m^{\mathrm{OPR}},t_{m+1}^{\mathrm{OPR}}],
\label{eq:time_to_angle}
\end{equation}
其中，$\theta_m^{\mathrm{OPR}}$ 为第 $m$ 个参考位置对应的累计转角；单槽 OPR 每转提供一次参考时，$\Delta\theta_m^{\mathrm{OPR}}=2\pi$，多槽或不等槽 OPR 则采用相邻参考位置之间的已知槽角。该映射利用实测 OPR 时刻反映不同转次或相邻槽区间的平均转速变化，并在相邻参考时刻之间采用角速度近似恒定的分段线性表示。

对第 $s$ 个传感器的第 $\ell$ 次目标叶片经过事件，记其名义参考转角为 $\theta_{s\ell}^{\mathrm{ref}}$，叶尖参考半径为 $R_{\mathrm{tip}}$。采样时刻 $t_{s\ell n}$ 对应的局部周向空间坐标定义为
\begin{equation}
x_{s\ell n}
=R_{\mathrm{tip}}
\left[\theta(t_{s\ell n})-\theta_{s\ell}^{\mathrm{ref}}\right].
\label{eq:local_space_coordinate}
\end{equation}
其中，$n$ 为单次叶片经过窗口内的采样点编号，$x=0$ 对应名义参考位置，正方向由转子旋转方向统一规定。式\eqref{eq:local_space_coordinate}把不同转速、不同转次下的时间采样点转换为具有相同物理含义的局部周向位置；若采用局部角坐标而非弧长坐标，则省略比例因子 $R_{\mathrm{tip}}$，但全文必须保持同一坐标单位。各传感器的局部周向坐标采用统一的物理单位、正方向和尺度，使后续共享的周向振动位移 $u(t)$ 在不同传感器模型中具有一致的物理含义。

设第 $s$ 个传感器在第 $\ell$ 个低速转次记录的原始电压为 $v_{s\ell}^{\mathrm{low}}(t_{s\ell n})$。将离散点对 $\{x_{s\ell n},v_{s\ell}^{\mathrm{low}}(t_{s\ell n})\}$ 插值到所有有效转次共同覆盖的空间网格 $\{x_p\}_{p=1}^{N_x}$，得到公共网格重采样波形
\begin{equation}
\overline v_{s\ell}^{\mathrm{low}}(x_p)
=\mathcal I
\left(
\{x_{s\ell n},v_{s\ell}^{\mathrm{low}}(t_{s\ell n})\};x_p
\right),
\label{eq:spatial_resampling}
\end{equation}
其中，$\mathcal I(\cdot)$ 表示空间插值算子。插值仅在各转次共同支持的有效空间区间内进行，不利用窗口外推值。剔除异常转次和无效采样点后，将 $L_s$ 个有效低速转次在同一空间位置上聚合，定义实测低速模板为
\begin{equation}
T_s^{\mathrm{low}}(x_p)
:=\operatorname{Aggregate}_{\ell=1,\ldots,L_s}
\left\{\overline v_{s\ell}^{\mathrm{low}}(x_p)\right\}.
\label{eq:low_template}
\end{equation}
其中，$L_s$ 为第 $s$ 个传感器的有效低速转次数，$\operatorname{Aggregate}$ 表示均值或预先规定的稳健聚合算子。$T_s^{\mathrm{low}}(x)$ 保存目标叶片、传感器安装、局部几何和测量链路共同形成的实际低速基准波形。

高速数据采用相同的式\eqref{eq:time_to_angle}和式\eqref{eq:local_space_coordinate}转换。后文分别以 $t_{sn}$、$x_{sn}$ 和 $v_{sn}$ 表示第 $s$ 个传感器第 $n$ 个高速观测点的采样时刻、局部周向坐标和实测电压。其中，$x_{sn}$ 用于查询空间波形，$t_{sn}$ 保留用于计算振动相位。由此，波形形状在空间域中比较，振动随时间的演化仍由原始采样时刻描述。

\subsection{静态间隙响应面与低速等效标定}

\subsubsection{间隙方向基函数与静态响应面构造}

静态标定用于建立无振动条件下电压关于径向间隙 $g$ 和局部周向坐标 $x$ 的响应关系。本文采用一个共享响应面作为统一间隙库，并假定不同传感器相对于该间隙库的主要差异可由后续的电压增益、偏置及空间配准参数表示；该假定应通过各传感器低速模板的重构误差检验。对固定的周向位置，本文采用常数项、反比项和对数项组成的低阶间隙基函数
\begin{equation}
\phi_0(g)=1,
\qquad
\phi_1(g)=\frac{1}{g},
\qquad
\phi_2(g)=\ln\left(\frac{g}{g_{\mathrm{ref}}}\right).
\label{eq:gap_basis}
\end{equation}
式中，$g_{\mathrm{ref}}>0$ 为参考间隙，用于使对数函数的自变量无量纲；$\phi_0$、$\phi_1$ 和 $\phi_2$ 分别为常数、反比和对数基函数。常数项给出间隙方向拟合的截距，反比项描述间隙减小时传感器响应增强的主要趋势，对数项修正有限标定范围内单一反比关系不能覆盖的缓变非线性。上述基函数是面向标定区间的经验表示，不作为严格的电场解析解。基于式\eqref{eq:gap_basis}，径向间隙--周向位置响应面定义为
\begin{equation}
\Rcal(g,x)
=\beta_0(x)+\frac{\beta_1(x)}{g}
+\beta_2(x)\ln\left(\frac{g}{g_{\mathrm{ref}}}\right),
\label{eq:response_surface}
\end{equation}
式中，$\Rcal(g,x)$ 为静态标定响应面；$\beta_0(x)$、$\beta_1(x)$ 和 $\beta_2(x)$ 为三个间隙基函数在周向位置 $x$ 处的系数，而不是预先规定形状的空间函数。允许这些系数随 $x$ 变化，是为了描述叶尖进入、正对和离开传感器敏感区域时，拟合截距、间隙灵敏度和非线性程度随周向位置发生的变化。

设静态标定包含 $N_g$ 个已知正间隙 $g_i$ 和 $N_x$ 个空间采样点 $x_p$。在固定空间位置 $x_p$ 处，将不同间隙下的实测电压和相应基函数值分别组成观测向量与标定矩阵：
\begin{equation}
\bm y_p=
\begin{bmatrix}
v_{1p}^{\mathrm{stat}} & \cdots & v_{N_gp}^{\mathrm{stat}}
\end{bmatrix}^{\mathrm T},
\qquad
\bm B_g=
\begin{bmatrix}
1 & 1/g_1 & \ln(g_1/g_{\mathrm{ref}})\\
\vdots & \vdots & \vdots\\
1 & 1/g_{N_g} & \ln(g_{N_g}/g_{\mathrm{ref}})
\end{bmatrix}.
\label{eq:cal_matrix}
\end{equation}
式中，$v_{ip}^{\mathrm{stat}}$ 为已知间隙 $g_i$、空间位置 $x_p$ 处的静态实测电压；$\bm y_p\in\mathbb R^{N_g}$ 为该位置的观测向量；$\bm B_g\in\mathbb R^{N_g\times3}$ 为间隙方向标定矩阵。三个系数可独立估计要求 $N_g\geq3$ 且 $\bm B_g$ 具有数值满列秩；若采用留一间隙重构检验，则应有 $N_g\geq4$。定义位置 $x_p$ 处的响应面系数向量为 $\bm\beta_p=[\beta_0(x_p),\beta_1(x_p),\beta_2(x_p)]^{\mathrm T}$，其最小二乘估计为
\begin{equation}
\bm\beta_p^*
=\arg\min_{\bm\beta_p}
\left\lVert \bm B_g\bm\beta_p-\bm y_p\right\rVert_2^2
\label{eq:cal_ls}
\end{equation}
式中，$\bm\beta_p^*$ 为 $\bm\beta_p$ 的估计值。具体而言，先在每个离散周向位置 $x_p$ 利用多个已知间隙估计三项系数，再沿周向位置对系数序列进行平滑插值，从而得到连续的 $\beta_j(x)$、响应面 $\Rcal(g,x)$ 及其间隙和空间导数。基函数形式及平滑结果通过留一间隙重构误差、$\partial\Rcal/\partial g$ 的符号与连续性以及标定矩阵的数值条件进行检验；后续识别仅在静态间隙库具有稳定覆盖的范围内进行。

\subsubsection{低速模板在静态响应面中的有界等效标定}

静态响应面提供统一的径向间隙--周向位置响应关系，不同传感器的安装零位、电压尺度和空间尺度通过传感器相关参数进行校正。为确定当前传感器在响应面中的低速查询路径，同时保留其实际基准波形，需要将实测模板 $T_s^{\mathrm{low}}(x)$ 与响应面进行有界匹配。

低速等效间隙路径和局部坐标到静态间隙库坐标的映射分别写为
\begin{equation}
g_s^{\mathrm{low}}(x)
=g_s^{(0)}+\mu_s^{\mathrm{cal}}(x-\tau_s),
\qquad
x_s^{\mathrm{lib}}(x)=\zeta_s(x-\tau_s).
\label{eq:low_cal_path}
\end{equation}
其中，$g_s^{(0)}$ 为低速模型等效基准径向间隙，$\mu_s^{\mathrm{cal}}$ 为径向间隙随局部周向位置变化的等效梯度，$\tau_s$ 为模板与间隙库的周向中心配准参数，$\zeta_s$ 为周向尺度修正系数。由此将低速查询路径明确写为静态响应面中的参数曲线
\begin{equation}
\Gamma_s(x)=
\left(g_s^{\mathrm{low}}(x),x_s^{\mathrm{lib}}(x)\right).
\label{eq:query_path}
\end{equation}
$\Gamma_s(x)$ 描述单个传感器有效窗口内，等效径向间隙随局部周向位置变化的标定路径。

为使响应面路径在有效标定域内逼近低速实测模板，定义路径参数向量 $\bm\vartheta_s=[g_s^{(0)},\mu_s^{\mathrm{cal}},\tau_s,\zeta_s]^{\mathrm T}$，并求解以下有界标定问题：
\begin{equation}
(\widehat{\bm\vartheta}_s,\widehat\kappa_s,\widehat b_s)
=\arg\min_{\substack{\bm\vartheta_s\in\Omega_s,\,\kappa_s>0\\b_s\in\mathbb R}}
\left\{
\sum_p
\left[
T_s^{\mathrm{low}}(x_p)-b_s
-\kappa_s\Rcal\!\left(
g_s^{\mathrm{low}}(x_p),x_s^{\mathrm{lib}}(x_p)
\right)
\right]^2
+P_{\mathrm{cal,out}}
\right\}
\label{eq:low_calibration}
\end{equation}
式中，$\Omega_s$ 为由静态间隙范围、有效空间窗口和物理可行范围共同确定的有界参数域，并在统一周向坐标正方向后约束 $\zeta_s>0$；$\kappa_s$ 和 $b_s$ 分别为响应面到第 $s$ 个传感器低速模板的电压增益和偏置；$P_{\mathrm{cal,out}}$ 为查询路径越出响应面支持范围时的惩罚项。统一电压极性后取 $\kappa_s>0$，以避免增益符号与响应方向相互补偿。对给定的 $\bm\vartheta_s$，$\kappa_s$ 和 $b_s$ 可由带上述符号约束的线性最小二乘确定，外层只需搜索查询路径参数。$b_s$ 仅为低速标定辅助参数，不进入高速主识别；$\kappa_s$ 用于将统一间隙库的电压变化尺度映射到当前传感器。

式\eqref{eq:low_calibration}用于确定低速模板在静态响应面中的等效查询路径，而不是用响应面拟合值替换实测模板。所得参数表示当前叶片与传感器组合在有效窗口内的等效标定关系。

低速标定结果通过增益方向、响应面覆盖范围以及重复数据和不同初值下查询路径的一致性进行检查。模板拟合残差与标定路径稳定性共同作为标定结果进入高速识别的依据。

\subsubsection{基于实测模板的间隙增量模型}

完成低速等效标定后，低速基准间隙和高速查询间隙分别定义为
\begin{align}
g_s^{\mathrm{base}}(x)
&=g_s^{(0)}+\mu_s^{\mathrm{cal}}(x-\tau_s),
\label{eq:g_base}\\
g_s^{\mathrm{qry}}(\Delta g_s,x)
&=g_s^{(0)}+\Delta g_s
+\mu_s^{\mathrm{cal}}(x-\tau_s).
\label{eq:g_query}
\end{align}
式中，$g_s^{\mathrm{base}}(x)$ 为低速标定路径上的基准间隙，$g_s^{\mathrm{qry}}(\Delta g_s,x)$ 为给定高速间隙增量后的查询间隙，$\Delta g_s$ 为第 $s$ 个传感器处高速状态相对于低速标定状态的等效径向间隙变化。两式使用同一个 $\mu_s^{\mathrm{cal}}$，因此高速径向间隙变化表示为低速标定路径沿 $g$ 方向的整体平移，同时保留径向间隙随局部周向位置变化的关系。

在相同空间查询位置 $x$ 处，将高速查询间隙与低速基准间隙对应的响应面输出相减，定义间隙变化引起的电压增量为
\begin{align}
\Delta\Rcal_s(\Delta g_s,x)
={}&\Rcal\!\left(
g_s^{\mathrm{qry}}(\Delta g_s,x),
x_s^{\mathrm{lib}}(x)
\right)
\nonumber\\
&-\Rcal\!\left(
g_s^{\mathrm{base}}(x),
x_s^{\mathrm{lib}}(x)
\right).
\label{eq:response_increment}
\end{align}
式中，$\Delta\Rcal_s(\Delta g_s,x)$ 表示静态响应面给出的相对电压变化。将该变化经传感器增益映射后叠加到实测低速模板上，定义传感器相关的间隙增量前向模型为
\begin{equation}
\Fsen_s(\Delta g_s,x)
=T_s^{\mathrm{low}}(x)
+\kappa_s\Delta\Rcal_s(\Delta g_s,x),
\label{eq:sensor_model}
\end{equation}
当 $\Delta g_s=0$ 时，式\eqref{eq:response_increment}严格为零，式\eqref{eq:sensor_model}退回实测模板 $T_s^{\mathrm{low}}(x)$。因此，低速模板保存绝对基准波形，静态响应面只提供相对低速状态的间隙增量响应，二者不会重复定义基准波形。

$T_s^{\mathrm{low}}(x)$、$g_s^{(0)}$、$\mu_s^{\mathrm{cal}}$、$\tau_s$、$\zeta_s$ 和 $\kappa_s$ 均在高速识别前确定，并在同一次高速主拟合中保持固定。重复低速标定和参数扰动用于评价这些标定量对高速间隙与振动结果的影响。

\subsection{高速完整波形联合模型}

\subsubsection{周向振动位移表示}

本文将振动识别量定义为叶尖在周向上的等效位移。当该位移可由有限个主导阶次及其邻近连续频率描述时，第 $s,n$ 个观测点的周向振动位移表示为
\begin{equation}
u(t_{sn})=\sum_{q=1}^{Q}
\left[
a_q\sin\psi_{q,sn}+b_q\cos\psi_{q,sn}
\right],
\label{eq:vibration}
\end{equation}
式中，$Q$ 为预设振动成分数，$a_q$ 和 $b_q$ 分别为第 $q$ 个成分的正弦和余弦系数，$\psi_{q,sn}$ 为该成分在观测点 $s,n$ 处的相位。为同时表示候选阶次和其邻近的连续频率偏移，$\psi_{q,sn}$ 定义为
\begin{equation}
\psi_{q,sn}
=k_q\theta_{sn}
+2\pi\Delta f_q(t_{sn}-t_c),
\label{eq:vibration_phase}
\end{equation}
式中，$k_q$ 为候选激振阶次，$\theta_{sn}=\theta(t_{sn})$ 为由式\eqref{eq:time_to_angle}得到的累计转角，$\Delta f_q$ 为相对阶次频率 $k_q\bar f_r$ 的连续频率偏移，$\bar f_r$ 为当前窗口的平均转频，$t_c$ 为窗口中心时刻。相应的物理频率、振动幅值和参考相位分别为
\begin{equation}
f_q=k_q\bar f_r+\Delta f_q,
\qquad
A_q=\sqrt{a_q^2+b_q^2},
\qquad
\varphi_q=\operatorname{atan2}(b_q,a_q).
\label{eq:amp_phase}
\end{equation}
式\eqref{eq:vibration}是多频一般形式。成分数 $Q$ 和候选阶次集合由分析频带、预设频谱范围或仿真设计给定，本文不把未知成分数判定并入同一优化问题。为避免成分标签交换及不同阶次--频率偏移组合重复表示同一频率，多频候选应满足
\begin{equation}
f_1<f_2<\cdots<f_Q,
\qquad
f_{q+1}-f_q\geq\Delta f_{\min}.
\label{eq:multifrequency_constraints}
\end{equation}
对候选阶次 $k$，若连续频率偏移的允许范围为 $[\Delta f_{k,\min},\Delta f_{k,\max}]$，则对应的物理频率搜索区间为
\begin{equation}
\left[
k\bar f_r+\Delta f_{k,\min},
k\bar f_r+\Delta f_{k,\max}
\right]
\label{eq:physical_frequency_interval}
\end{equation}
式中，$\Delta f_{k,\min}$ 和 $\Delta f_{k,\max}$ 分别为阶次 $k$ 对应的连续频率偏移下限和上限。不同候选阶次对应的物理频率区间应互不重叠。$\Delta f_{\min}$ 是由分析窗口长度和仿真分辨性试验确定的最低分离尺度，不是普适常数。当变量投影矩阵数值病态或近频列无法稳定分离时，应将相关成分标记为不可分辨，而不强行报告两个独立可信的幅值和相位。本文试验分析采用 $Q=1$，在给定搜索频带内比较候选阶次；$Q>1$ 仅作为一般模型并留待仿真验证。

\subsubsection{高速波形观测方程}

按照第2节的建模假设，高速相对于低速的主要静态变化表示为径向间隙增量，叶片振动表示为周向等效位移。由于各传感器的局部坐标已在第3.2节统一，全部传感器共享同一周向振动位移函数 $u(t)$。

叶片振动和空间对准误差通过改变模板取值位置进入前向模型。第 $s,n$ 个观测点的有效周向坐标定义为
\begin{equation}
\xi_{sn}=x_{sn}-d-u(t_{sn}),
\label{eq:xi}
\end{equation}
其中，$d$ 为当前数据窗口内全部传感器共享的空间对准偏移，用于表征并补偿同一坐标构造过程产生的公共零位偏差。各传感器的固定安装角差已通过式\eqref{eq:local_space_coordinate}中的名义参考转角 $\theta_{s\ell}^{\mathrm{ref}}$ 纳入局部坐标定义，因此高速主模型不再引入逐传感器空间偏移。本文规定 $u(t)>0$ 表示叶尖沿局部周向坐标正方向移动；在固定名义观测坐标下，模板与响应面的有效取值位置相应减小，故式\eqref{eq:xi}中振动位移前取负号。

将式\eqref{eq:xi}代入式\eqref{eq:sensor_model}，得到固定低速标定梯度模型的完整波形预测：
\begin{equation}
\widehat v_{sn}(\bm\theta_c;\bm k)
=\Fsen_s\!\left(
\Delta g_s,
x_{sn}-d-u(t_{sn})
\right).
\label{eq:observation}
\end{equation}
式\eqref{eq:observation}是本文的核心前向模型。高速间隙增量 $\Delta g_s$ 改变响应面的间隙输入，公共空间偏移和振动位移改变模板及响应面的空间取值位置，固定的 $\mu_s^{\mathrm{cal}}$ 则在该有效周向位置处参与间隙路径计算。因此，已标定梯度与振动位移之间的几何作用被保留，但梯度本身不作为高速自由参数。上述效应通过同一个非线性前向模型共同形成预测波形，不能解释为彼此独立的加性电压项。

\subsubsection{识别参数与目标函数}

将候选阶次组成离散向量 $\bm k=[k_1,\ldots,k_Q]^{\mathrm T}$，其有限候选集合记为 $\mathcal K$。给定 $\bm k$ 后，连续待估参数写为
\begin{equation}
\bm\theta_c=
\left[
\Delta g_1,\ldots,\Delta g_{N_s},d,
a_1,b_1,\Delta f_1,\ldots,
a_Q,b_Q,\Delta f_Q
\right]^{\mathrm T}.
\label{eq:theta}
\end{equation}
其中不包含高速自由梯度增量。记具有有效间隙标定信息的传感器集合为 $\mathcal S_g$；对 $s\notin\mathcal S_g$，在整个高速主优化中约束 $\Delta g_s\equiv0$，但相应观测仍参与公共偏移和振动参数估计。为区分候选生成和最终参数估计，变量投影阶段使用预先确定的对角权重矩阵 $\bm\Lambda_{\mathrm{VP}}$，以减小低灵敏度点和响应面边缘点对候选生成的影响；最终候选比较采用未加权完整电压残差，并加入响应面越界、无效查询和间隙增量正则项。定义
\begin{equation}
\Phi(\bm\theta_c;\bm k)=
\sum_{s,n}
\left[v_{sn}-\widehat v_{sn}(\bm\theta_c;\bm k)\right]^2
+P_{\mathrm{out}}(\bm\theta_c;\bm k)
+P_{\mathrm{invalid}}(\bm\theta_c;\bm k)
+P_g(\Delta\bm g).
\label{eq:objective}
\end{equation}
$P_{\mathrm{out}}$ 用于抑制查询点超出静态响应面和低速模板的有效范围，$P_{\mathrm{invalid}}$ 用于惩罚无效插值，$P_g$ 用于限制缺少独立间隙真值时的过度静态修正。令 $o_{sn}$ 为第 $s,n$ 个查询点到有效标定域的越界距离，$N_{\mathrm{invalid}}$ 为无效查询点数，$N$ 为当前窗口观测点数，可采用
\begin{equation}
\begin{aligned}
P_{\mathrm{out}}&=\gamma_{\mathrm{out}}^2\sum_{s,n}o_{sn}^2,\\
P_{\mathrm{invalid}}&=M_{\mathrm{invalid}}N_{\mathrm{invalid}},\\
P_g&=
\frac{N\gamma_g^2}{N_s}
\sum_{s=1}^{N_s}
\left(\frac{\Delta g_s}{g_{s,\mathrm{scale}}}\right)^2,
\end{aligned}
\label{eq:penalties}
\end{equation}
其中，间隙坐标和周向坐标均换算为长度单位后，$o_{sn}$ 取查询间隙和周向位置相对于各自有效区间的二维欧氏越界量；$\gamma_{\mathrm{out}}$ 的量纲用于将该越界量换算为电压残差尺度。$M_{\mathrm{invalid}}$ 的量纲为电压平方，$\gamma_g$ 为间隙增量正则项的电压尺度；这些系数应在全部数据集上保持固定。$g_{s,\mathrm{scale}}$ 为由标定范围或独立物理信息给出的间隙尺度。不同方法比较时应保持相同观测点、完整波形残差定义和惩罚口径。式\eqref{eq:objective}中的未加权残差是最终结果口径，不与变量投影阶段的加权残差混用。该口径默认各传感器的有效采样点数和电压尺度具有可比性；否则，总残差会自然偏向点数较多或电压尺度较大的传感器。本文保留未加权总残差作为统一优化准则，同时通过逐传感器误差和传感器组合敏感性检查结果是否受单个传感器主导。

为明确模型层级，固定模板模型对应 $\Delta g_s=0$；固定低速标定梯度模型允许 $\Delta g_s$ 变化但固定 $\mu_s^{\mathrm{cal}}$；允许高速自由梯度增量的模型仅作为敏感性对照，不作为主振动幅值模型。

\subsection{投影局部搜索与完整模型优化}
\label{sec:candidate_solution}

直接从任意初值同时搜索阶次、连续频率、间隙、公共空间偏移和振动参数容易进入局部极小。本文不设置忽略振动的独立静态初值步骤，而是对每组候选阶次及连续频率参数先计算振动候选，再在该候选下构造间隙搜索范围，最后对全部保留候选进行完整非线性优化。

\subsubsection{加权变量投影候选生成}

对给定候选阶次向量 $\bm k$、连续频率偏移向量 $\Delta\bm f$、参考间隙增量 $\Delta\bm g^{\mathrm{ref}}$ 和参考公共空间偏移 $d^{\mathrm{ref}}$，先在参考状态附近对振动位移作一阶展开。$d^{\mathrm{ref}}$ 在预设的有界公共偏移区间内作为外层候选参数；本文取 $\Delta\bm g^{\mathrm{ref}}=\bm 0$，该零值只是投影中心构造所采用的参考状态，不是高速间隙不变的物理假设。参考状态的无振动预测值及其总空间导数定义为
\begin{equation}
F_{sn}^{(0)}
=\Fsen_s(\Delta g_s^{\mathrm{ref}},x_{sn}-d^{\mathrm{ref}}),
\qquad
F_{x,sn}^{(0)}=
\left.\frac{\partial\Fsen_s(\Delta g_s,x)}{\partial x}
\right|_{\substack{\Delta g_s=\Delta g_s^{\mathrm{ref}}\\
x=x_{sn}-d^{\mathrm{ref}}}},
\label{eq:local_derivative}
\end{equation}
式中，$F_{sn}^{(0)}$ 为参考间隙和参考公共偏移下的无振动预测电压，$F_{x,sn}^{(0)}$ 为固定低速标定梯度前向模型对空间查询坐标的总导数。由此，观测点 $s,n$ 处的波形关于周向振动位移的一阶近似为
\begin{equation}
\widehat v_{sn}
\approx F_{sn}^{(0)}-F_{x,sn}^{(0)}u(t_{sn}).
\label{eq:linearized_wave}
\end{equation}
$F_{x,sn}^{(0)}$ 包含低速模板形状、坐标映射以及已标定梯度对空间查询的共同影响。

对给定阶次向量 $\bm k$ 和频率偏移向量 $\Delta\bm f$，式\eqref{eq:linearized_wave}对 $a_q,b_q$ 为线性关系。将全部观测点按固定顺序堆叠，得到线性候选模型
\begin{equation}
\bm r_0(d^{\mathrm{ref}};\Delta\bm g^{\mathrm{ref}})
\approx
\bm D(\bm k,\Delta\bm f,d^{\mathrm{ref}};\Delta\bm g^{\mathrm{ref}})
\bm\alpha+\bm e,
\qquad
\bm\alpha=
[a_1,b_1,\ldots,a_Q,b_Q]^{\mathrm T},
\label{eq:vp_model}
\end{equation}
式中，$\bm v$ 和 $\bm F^{(0)}$ 分别为实测电压与参考预测电压的堆叠向量，$\bm r_0=\bm v-\bm F^{(0)}$ 为参考残差，$\bm\alpha$ 为全部振动线性系数组成的向量，$\bm D$ 为相应设计矩阵，$\bm e$ 为局部线性化余项。矩阵 $\bm D$ 中对应第 $q$ 个正弦和余弦系数的两列定义为
\begin{equation}
D_{sn,2q-1}=-F_{x,sn}^{(0)}\sin\psi_{q,sn},
\qquad
D_{sn,2q}=-F_{x,sn}^{(0)}\cos\psi_{q,sn}.
\label{eq:vp_design}
\end{equation}
变量投影阶段使用预先确定的对角权重矩阵 $\bm\Lambda_{\mathrm{VP}}$，降低低空间灵敏度点和响应面边缘点对候选生成的影响。为使投影和范数均在普通欧氏空间中解释，白化残差和设计矩阵定义为
\begin{equation}
\widetilde{\bm r}_0=\bm\Lambda_{\mathrm{VP}}^{1/2}\bm r_0,
\qquad
\widetilde{\bm D}=\bm\Lambda_{\mathrm{VP}}^{1/2}\bm D.
\label{eq:vp_whitening}
\end{equation}
在每组候选阶次、连续频率偏移和公共空间偏移参考值下，振动线性系数由变量投影求得\cite{Golub1973,Golub2003}：
\begin{equation}
\widehat{\bm\alpha}
(\bm k,\Delta\bm f,d^{\mathrm{ref}};\Delta\bm g^{\mathrm{ref}})
=\widetilde{\bm D}
(\bm k,\Delta\bm f,d^{\mathrm{ref}};\Delta\bm g^{\mathrm{ref}})^{\dagger}
\widetilde{\bm r}_0
(d^{\mathrm{ref}};\Delta\bm g^{\mathrm{ref}})
\label{eq:vp_solution}
\end{equation}
式中，$(\cdot)^{\dagger}$ 表示通过奇异值分解计算的 Moore--Penrose 广义逆。将线性系数代回白化局部模型后，变量投影候选得分定义为
\begin{equation}
\Psi_{\mathrm{VP}}
(\bm k,\Delta\bm f,d^{\mathrm{ref}};\Delta\bm g^{\mathrm{ref}})
=\left\lVert
\widetilde{\bm r}_0-
\widetilde{\bm D}\widehat{\bm\alpha}
\right\rVert_2^2.
\label{eq:vp_score}
\end{equation}
外层按式\eqref{eq:vp_score}比较候选阶次、连续频率偏移和公共空间偏移参考值，内层直接求解 $a_q,b_q$，从而减少候选生成阶段的非线性搜索维数。为保证线性振动系数在当前候选下可数值分离，$\widetilde{\bm D}$ 经统一奇异值阈值截断后的数值秩应为 $2Q$，且最小保留奇异值或条件数满足预设稳定性要求；不满足该条件的组合不进入保留候选集合 $\mathcal C(\bm k)$。

式\eqref{eq:vp_solution}只用于为每组外层候选求得相应的振动线性系数；与该组系数配对的 $d^{\mathrm{ref}}$ 记为 $d^{\mathrm{cand}}$。在满足数值秩条件的组合中，保留投影得分较小的有限候选。具体保留数量或得分阈值属于统一设置的求解参数，不改变后续完整模型的判定准则。这些量均为完整优化的候选初值，不作为最终振动或空间偏移结果。所有保留候选仍需进入式\eqref{eq:observation}的完整非线性模型，不能用一阶近似残差替代最终完整波形残差。

\subsubsection{条件间隙灵敏度与局部搜索范围}

对每个保留候选，将变量投影得到的振动系数与对应的公共空间偏移、频率参数和参考间隙增量组合为
\begin{equation}
\bm\theta^{\mathrm{cand}}=
\left\{
\Delta\bm g^{\mathrm{ref}},
d^{\mathrm{cand}},
\bm k^{\mathrm{cand}},
\widehat{\bm a},
\widehat{\bm b},
\Delta\bm f^{\mathrm{cand}}
\right\},
\label{eq:candidate_point}
\end{equation}
其中，$\widehat{\bm a}$ 和 $\widehat{\bm b}$ 由式\eqref{eq:vp_solution}中的 $\widehat{\bm\alpha}$ 分解得到。在该候选点代入完整非线性前向模型，定义
\begin{equation}
\bm r^{\mathrm{cand}}
=\bm v-\widehat{\bm v}(\bm\theta^{\mathrm{cand}}).
\label{eq:candidate_residual}
\end{equation}
下文的 $\bm J_g$ 和 $\bm J_{\nu}$ 均定义为预测波形 $\widehat{\bm v}$ 对相应参数的雅可比，并在 $\bm\theta^{\mathrm{cand}}$ 处由完整模型求导或差分得到，即
\begin{equation}
\bm J_g=
\left.
\frac{\partial\widehat{\bm v}}
{\partial\Delta\bm g}
\right|_{\bm\theta^{\mathrm{cand}}},
\qquad
\bm J_{\nu}=
\left.
\frac{\partial\widehat{\bm v}}
{\partial\bm\nu}
\right|_{\bm\theta^{\mathrm{cand}}}.
\label{eq:prediction_jacobians}
\end{equation}
离散阶次 $\bm k^{\mathrm{cand}}$ 在局部投影中保持固定，不构造对阶次的微分灵敏度。为简化记号，以下将 $\bm r^{\mathrm{cand}}$ 记为 $\bm r$。当参数修正量为 $\delta\bm g$ 和 $\delta\bm\nu$ 时，修正后的残差在该候选点附近满足
\begin{equation}
\bm r_{\mathrm{new}}
\approx
\bm r-
\bm J_g\delta\bm g-
\bm J_{\nu}\delta\bm\nu,
\label{eq:local_model}
\end{equation}
其中，$\bm J_g$ 由具有间隙标定信息的传感器集合 $\mathcal S_g$ 对应的间隙灵敏度列组成。干扰参数灵敏度矩阵写为
\begin{equation}
\bm J_{\nu}=
[\bm J_d,\bm J_{a},\bm J_b,\bm J_{\Delta f}],
\label{eq:nuisance_jacobian}
\end{equation}
其各列分别对应公共空间偏移、全部振动正弦和余弦系数以及连续频率偏移。$\delta\bm g$ 和 $\delta\bm\nu$ 为参考状态附近的局部参数修正量。未进入 $\mathcal S_g$ 的传感器仍参与公共偏移、振动和频率方向的估计，但其间隙增量固定为零。

采用与变量投影相同的权重进行白化：
\begin{equation}
\widetilde{\bm r}=\bm\Lambda_{\mathrm{VP}}^{1/2}\bm r,
\qquad
\widetilde{\bm J}_{g}=\bm\Lambda_{\mathrm{VP}}^{1/2}\bm J_g,
\qquad
\widetilde{\bm J}_{\nu}=\bm\Lambda_{\mathrm{VP}}^{1/2}\bm J_{\nu}.
\label{eq:gap_whitening}
\end{equation}
通过奇异值分解得到 $\widetilde{\bm J}_{\nu}$ 的数值列空间正交基 $\bm U_{\nu}$，并定义
\begin{equation}
\bm P_{\nu}=\bm U_{\nu}\bm U_{\nu}^{\mathrm T},
\qquad
\bm Q_{\nu}=\bm I-\bm P_{\nu}.
\label{eq:projector}
\end{equation}
$\bm P_{\nu}$ 投影到公共偏移、振动和频率灵敏度张成的局部子空间，$\bm Q_{\nu}$ 则保留这些方向不能一阶解释的部分。由此得到联合间隙矩阵和投影残差：
\begin{equation}
\bm G=\bm Q_{\nu}\widetilde{\bm J}_g,
\qquad
\bm r_{\perp}=\bm Q_{\nu}\widetilde{\bm r}.
\label{eq:projected_gap_system}
\end{equation}
对 $\bm G$ 进行奇异值分解，并仅保留大于既定相对阈值的奇异值。记由此得到的截断广义逆为 $\bm G_{\mathrm{tr}}^{\dagger}$。只有当保留的数值秩等于待估间隙参数个数时，才计算联合局部线性间隙修正量
\begin{equation}
\delta\bm g^{\mathrm{LS}}
=\bm G_{\mathrm{tr}}^{\dagger}\bm r_{\perp}.
\label{eq:joint_gap_ls}
\end{equation}
若 $\bm G$ 的数值秩不足，则不计算联合投影搜索中心；对全部 $s\in\mathcal S_g$ 取 $\overline{\Delta g}_s=\Delta g_s^{\mathrm{ref}}$，并采用预设的保守半宽 $h_{g,s}^{\mathrm{cons}}$。奇异值截断阈值在全部数据集上保持一致，避免因候选而改变秩判据。
为评价第 $s$ 个间隙参数，在已去除干扰参数方向的基础上，还需排除其他间隙参数的一阶灵敏度方向。记 $\bm G_{-s}$ 为删去第 $s$ 列后的 $\bm G$，$\bm P_{g,-s}$ 为通过奇异值分解得到的 $\bm G_{-s}$ 数值列空间正交投影，$\bm g_s$ 为 $\bm G$ 的第 $s$ 列，$\widetilde{\bm j}_{g,s}$ 为白化间隙雅可比矩阵 $\widetilde{\bm J}_g$ 的第 $s$ 列；若不存在其他有效间隙参数，则取 $\bm P_{g,-s}=\bm0$。由此定义
\begin{equation}
\bm q_{g,s}=
\left(\bm I-\bm P_{g,-s}\right)\bm g_s,
\qquad
\eta_{g,s}=
\frac{\left\lVert\bm q_{g,s}\right\rVert_2}
{\left\lVert\widetilde{\bm j}_{g,s}\right\rVert_2}.
\label{eq:gap_separability}
\end{equation}
$\eta_{g,s}$ 描述第 $s$ 个间隙参数的一阶灵敏度中，不能被公共空间偏移、振动、频率偏移及其他间隙参数灵敏度方向解释的比例。因此，$\eta_{g,s}$ 是当前候选附近的局部条件可辨识性指标；其值较小表示该参数的一阶灵敏度接近混叠，但它本身不是间隙误差上界。即使原始间隙灵敏度列在不同传感器观测区间上互不重叠，经公共干扰子空间投影后也不再假定其严格正交，因此仍需按联合间隙雅可比矩阵进行条件化处理。

相应的条件投影残差为
\begin{equation}
\bm r_{\perp,s}
=\left(\bm I-\bm P_{g,-s}\right)\bm r_{\perp}.
\label{eq:conditional_gap_residual}
\end{equation}
条件化残差与条件灵敏度之比用于定义局部搜索尺度基准：
\begin{equation}
 h_{g,s}^{\mathrm{proj}}
:=\frac{\left\lVert\bm r_{\perp,s}\right\rVert_2}
{\left\lVert\bm q_{g,s}\right\rVert_2}.
\label{eq:gap_bound}
\end{equation}
该比值反映在当前局部线性模型中，未解释残差相对于独立间隙灵敏度的尺度；它不直接作为间隙估计值。利用联合最小二乘得到的有符号修正量更新局部搜索中心：
\begin{equation}
\overline{\Delta g}_s
=\operatorname{clip}\!\left(
\Delta g_s^{\mathrm{ref}}
+\delta g_s^{\mathrm{LS}},
\Delta g_{s,\min}^{\mathrm{adm}},
\Delta g_{s,\max}^{\mathrm{adm}}
\right),
\label{eq:gap_center}
\end{equation}
其中，$\operatorname{clip}$ 表示按传感器分量投影到各自允许区间；$\Delta g_{s,\min}^{\mathrm{adm}}$ 和 $\Delta g_{s,\max}^{\mathrm{adm}}$ 由响应面覆盖范围与物理允许范围共同确定。本文取 $\Delta g_s^{\mathrm{ref}}=0$，因此投影联合解直接构成候选相关的间隙搜索中心。实际局部搜索半宽设置为
\begin{equation}
h_{g,s}
=\operatorname{clip}\!\left(
\rho_g h_{g,s}^{\mathrm{proj}},
h_{g,\min},
h_{g,\max}
\right),
\label{eq:local_range}
\end{equation}
其中，$\rho_g>1$ 为保守放大系数，$h_{g,\min}$ 和 $h_{g,\max}$ 分别避免搜索范围过窄或超出间隙库支持范围。$h_{g,s}^{\mathrm{proj}}$ 是投影残差驱动的局部搜索尺度，不是统计置信区间，也不是非线性问题中真实间隙误差的严格上界。除上述联合矩阵秩不足时的整体回退外，若仅第 $s$ 个参数的 $\eta_{g,s}$ 或 $\|\bm q_{g,s}\|_2$ 低于阈值，则只对该传感器取 $\overline{\Delta g}_s=\Delta g_s^{\mathrm{ref}}$ 和 $h_{g,s}=h_{g,s}^{\mathrm{cons}}$；其余满足条件的传感器仍采用投影中心和半宽。若干扰子空间对奇异值截断阈值过度敏感，则按联合秩不足处理。若最终结果命中允许边界，也应标记为局部弱可辨识，不能据此声称全局唯一性。

\subsubsection{完整模型有界优化与候选选择}

对每个保留的候选，以上述投影中心和半宽构造间隙边界，并在候选附近对完整前向模型进行有界联合优化。主要局部范围写为
\begin{equation}
\begin{aligned}
|\Delta g_s-\overline{\Delta g}_s|&\leq h_{g,s},\\
|d-d^{\mathrm{cand}}|&\leq h_d,\\
|\Delta f_q-\Delta f_q^{\mathrm{cand}}|&\leq h_{f,q}.
\end{aligned}
\label{eq:search_domain}
\end{equation}
$d^{\mathrm{cand}}$ 和 $\Delta f_q^{\mathrm{cand}}$ 分别为变量投影候选给出的公共偏移和连续频率参考值。振动系数、间隙增量、公共偏移和连续频率均在式\eqref{eq:observation}的完整非线性模型中更新。加权残差只用于候选生成和投影搜索范围构造；每个保留候选的局部优化和最终排序均采用式\eqref{eq:objective}规定的同一未加权完整目标，包括响应面越界、无效查询和间隙增量正则项，从而保证候选局部优化准则与最终排序准则一致。

对给定离散阶次向量 $\bm k$，记保留候选集合为 $\mathcal C(\bm k)$，由第 $j$ 个候选及式\eqref{eq:search_domain}构成的局部有界域为 $\Omega_j(\bm k)$。各候选的完整模型局部解及最终候选选择分别写为
\begin{equation}
\widehat{\bm\theta}_{c,j}(\bm k)
=\arg\min_{\bm\theta_c\in\Omega_j(\bm k)}
\Phi(\bm\theta_c;\bm k),
\qquad
(\widehat{\bm k},\widehat j)
=\arg\min_{\substack{\bm k\in\mathcal K\\j\in\mathcal C(\bm k)}}
\Phi\!\left(\widehat{\bm\theta}_{c,j}(\bm k);\bm k\right).
\label{eq:discrete_continuous_objective}
\end{equation}
$\Omega_j(\bm k)$ 同时受总体标定范围、物理边界和候选局部范围约束，最终主参数记为 $\widehat{\bm\theta}=\{\widehat{\bm\theta}_{c,\widehat j}(\widehat{\bm k}),\widehat{\bm k}\}$。式\eqref{eq:discrete_continuous_objective}表示有限候选之间的局部优化与选择，不声称由单一初值获得连续参数问题的全局极小值。

最终输出为
\begin{equation}
\left\{
\widehat{\Delta g}_s,
\widehat d,
\widehat k_q,
\widehat{\Delta f}_q,
\widehat A_q,
\widehat\varphi_q
\right\},
\label{eq:outputs}
\end{equation}
其中，$\widehat f_q=\widehat k_q\bar f_r+\widehat{\Delta f}_q$，$\widehat A_q$ 和 $\widehat\varphi_q$ 由式\eqref{eq:amp_phase}计算。若多个传感器在物理上应共享同一个高速间隙变化，可进一步施加 $\Delta g_s=\Delta g$ 或跨传感器正则化约束；该约束应由传感器布置和被测结构决定，不能仅因其降低残差而采用。

\section{仿真验证方案}

仿真验证围绕径向间隙变化能否从周向振动位移中分离展开。除下列四类主工况外，可在各工况中加入响应面扰动、低速标定参数扰动、OPR 时间误差、空间坐标配准误差、异常点和未建模弱频率成分，以评价方法在非理想条件下的识别稳定性。

第一类为基线工况。令高速径向间隙增量为零，且标定参数与数据生成条件一致，用于检查各方法对周向振动频率、幅值和相位的基本识别能力。

第二类为高速径向间隙变化工况。在保持低速标定关系正确的条件下，分别施加全传感器共同间隙变化和单传感器等效间隙变化，比较固定模板模型与固定低速标定梯度模型的间隙误差、振动幅值误差、位移均方根误差和完整波形残差。该组仿真用于检验径向间隙变化对周向振动识别的影响及本文方法的分离效果。

第三类为低速标定扰动工况。对低速等效局部间隙梯度和空间配准参数施加不同幅值的标定扰动，报告径向间隙增量和周向振动参数的识别误差，用于评价低速标定重复性对高速结果的影响。

第四类为多频工况。设置频率间隔、幅值比和相位差不同的 $Q>1$ 振动成分，报告各成分的阶次、连续频率、幅值和相位误差，以及变量投影矩阵的条件数。该组仿真用于确定有限窗口内可稳定分离的最小频率间隔和幅值范围。

比较方法至少包括固定模板模型和固定低速标定梯度模型。求解策略消融比较固定间隙范围、候选相关的投影局部搜索和直接全参数非线性优化，报告候选覆盖率、函数评价次数、边界命中率、收敛成功率和最终误差。全部有效窗口采用相同的搜索中心与搜索半宽构造方法、相同的数值阈值及退化回退规则，并报告保守回退比例和边界命中情况。

\section{试验验证方法}

试验标定分为两个步骤。第一步，由满足低振动和重复性条件的低速旋转数据经公共空间网格重采样和多转聚合得到实测模板 $T_s^{\mathrm{low}}(x)$；第二步，将该模板与预先建立的静态多间隙响应面按式\eqref{eq:low_calibration}进行有界匹配，得到 $g_s^{(0)}$、$\mu_s^{\mathrm{cal}}$、$\tau_s$、$\zeta_s$ 和 $\kappa_s$。因此，上述标定参数由低速模板和静态间隙库共同确定。试验应明确报告所采用的聚合算子、异常转次剔除规则、有效低速转次数、模板重复性、标定边界、不同初值的一致性和重复标定误差；高速识别阶段保持 $\mu_s^{\mathrm{cal}}$ 固定。

主对比应包括固定模板模型和固定低速标定梯度模型。报告量包括识别阶次、连续频率、振动幅值、总未加权完整波形残差、$\widehat{\Delta g}_s$、$\eta_{g,s}$、$h_{g,s}^{\mathrm{proj}}$、$h_{g,s}$、边界命中情况以及结果随数据窗口的重复性。为揭示采样点数量和电压尺度差异对总残差构成的影响，并避免总残差指标掩盖单个传感器的拟合差异，还应逐传感器报告
\begin{equation}
\mathrm{RMSE}_s=
\sqrt{\frac{1}{N_{\mathrm{obs},s}}\sum_{n=1}^{N_{\mathrm{obs},s}}
\left(v_{sn}-\widehat v_{sn}\right)^2},
\label{eq:sensor_rmse}
\end{equation}
其中 $N_{\mathrm{obs},s}$ 为第 $s$ 个传感器纳入拟合的有效观测点数。若不同传感器组合得到不同振动幅值，还应分别检查间隙增量、低速等效标定误差和单传感器残差对结果的影响。

若有应变片、激光测振或其他独立振动参考，其作用是验证频率、阶次或幅值。若没有独立间隙真值，则 $\widehat{\Delta g}_s$ 解释为相对于低速基准的模型等效径向间隙变化，不据此评价绝对间隙测量精度。

\section{讨论}

固定低速标定梯度模型的主要作用是把高速识别前可获得的信息与需要由高速波形估计的信息分开。低速旋转数据用于建立传感器相关的实测基准模板；静态多间隙响应面提供径向间隙方向的标定关系；二者通过有界匹配共同确定低速等效查询路径。高速数据用于估计相对径向间隙变化、公共空间偏移和周向振动。该划分使径向间隙变化引起的脉冲幅值和形状变化与周向振动引起的脉冲位置变化在同一前向模型中得到描述。

等效局部间隙梯度描述径向间隙随叶尖局部周向位置的变化关系。高速识别保持这一低速标定关系不变，使周向振动改变查询位置时，相应的局部等效间隙变化仍由同一几何路径计算。低速重复标定和参数扰动试验用于评价该标定关系的稳定性。

各传感器的 $\Delta g_s$ 用于描述相应测点处相对于低速状态的径向间隙变化。若传感器布置和被测结构决定不同测点具有共同的间隙变化，则可采用共享间隙或跨传感器约束，并通过独立间隙参考进行验证。

本文模型在静态响应面具有稳定覆盖、低速模板具有良好重复性且主要振动成分位于候选频率范围内使用。局部条件可辨识性指标、边界命中情况和逐传感器残差共同用于评价识别结果。

\section{结论}

径向间隙变化会改变叶尖脉冲的幅值和形状，并影响传统到达时刻提取及完整波形匹配得到的周向位移。本文先由低速旋转数据建立实测模板，再利用静态多间隙响应面对模板进行有界等效标定，得到基准径向间隙、等效局部间隙梯度、空间配准和增益参数；高速阶段联合估计相对径向间隙增量、公共空间对准偏移和周向振动参数。低速标定梯度通过振动修正后的有效周向坐标进入完整前向模型，从而保留局部几何与周向振动之间的作用关系。

方法使用加权局部线性化变量投影产生候选，并利用条件化间隙灵敏度、有符号局部修正量和投影残差构造候选相关的间隙搜索中心与搜索半宽。该半宽是数值局部搜索尺度，不是统计置信区间。保留候选均在未加权完整非线性目标下进行有界联合优化和排序。

该模型识别相对于低速标定状态的高速等效径向间隙变化，并给出叶尖周向振动的频率、幅值和相位。其应用以低速模板重复性、模板与静态响应面的等效标定稳定性以及响应面覆盖范围为基础。本文试验分析采用单主导频率特例，多频形式仍需通过仿真验证；缺少独立间隙真值时，高速间隙结果表述为相对于低速状态的模型等效变化量。

\section*{主要符号}

\begin{longtable}{>{\centering\arraybackslash}p{0.24\textwidth}p{0.68\textwidth}}
\toprule
符号 & 含义\\
\midrule
\endfirsthead
\toprule
符号 & 含义\\
\midrule
\endhead
$s,n,q$ & 传感器、观测点和振动成分编号\\
$N_s,Q$ & 参与识别的传感器数和预设振动成分数\\
$L_s,N_{\mathrm{obs},s}$ & 第 $s$ 个传感器的有效低速转次数和高速有效观测点数\\
$v_{s\ell}^{\mathrm{low}}(t),\overline v_{s\ell}^{\mathrm{low}}(x)$ & 第 $s$ 个传感器第 $\ell$ 个低速转次的时间域原始波形和公共空间网格重采样波形\\
$v_{sn},\widehat v_{sn}$ & 实测电压和完整模型预测电压\\
$V_{\mathrm{tip}},\Delta t_u,\delta t_g$ & 叶尖周向速度、振动引起的到达时刻偏差和间隙变化引起的特征定位偏差\\
$t_m^{\mathrm{OPR}},\theta(t),\Delta\theta_m^{\mathrm{OPR}}$ & OPR参考时刻、累计转角和相邻参考位置之间的已知转角增量\\
$\theta_{s\ell}^{\mathrm{ref}},R_{\mathrm{tip}},x_{s\ell n}$ & 目标叶片经过的名义参考转角、叶尖参考半径和局部周向空间坐标\\
$t_{sn},x_{sn},\xi_{sn}$ & 采样时刻、名义局部周向坐标和振动修正后的有效周向坐标\\
$g,x$ & 等效径向间隙和局部周向坐标\\
$T_x(x),S_g(x)$ & 问题分析中低速模板的空间导数和基准间隙路径附近的间隙灵敏度\\
$\phi_j(g)$ & 静态响应面的间隙方向基函数\\
$\Rcal(g,x)$ & 径向间隙--周向位置静态标定响应面\\
$\Fsen_s(\Delta g_s,x)$ & 固定低速标定梯度的传感器相关前向模型\\
$\Gamma_s(x)$ & 低速模板在静态响应面中的等效径向间隙--周向位置查询路径\\
$T_s^{\mathrm{low}}(x)$ & 第 $s$ 个传感器由低速重复波形聚合得到的实测基准模板\\
$g_s^{(0)}$ & 第 $s$ 个传感器的低速等效基准径向间隙\\
$\Delta g_s$ & 高速状态相对低速基准的径向间隙增量\\
$\mu_s^{\mathrm{cal}}$ & 径向间隙随局部周向位置变化的低速等效标定梯度\\
$\tau_s,\zeta_s,\kappa_s$ & 模板--间隙库中心配准、空间尺度修正和电压增益修正参数\\
$b_s,\Omega_s,P_{\mathrm{cal,out}}$ & 低速标定辅助偏置、有界参数域和标定越界惩罚\\
$d$ & 当前数据窗口内全部传感器共享的空间对准偏移\\
$u(t)$ & 叶尖振动在周向上的等效位移\\
$k_q,\Delta f_q,\psi_{q,sn}$ & 候选阶次、连续频率偏移和振动相位\\
$\Delta f_{\min}$ & 多频成分在当前分析条件下允许独立报告的最低频率间隔\\
$a_q,b_q$ & 振动正弦、余弦系数\\
$A_q,\varphi_q$ & 振动幅值和参考相位\\
$\lambda_{sn}^{\mathrm{VP}},\bm\Lambda_{\mathrm{VP}}$ & 变量投影阶段的观测点权重和对角权重矩阵\\
$\Psi_{\mathrm{VP}},\mathcal C(\bm k),\Omega_j(\bm k)$ & 变量投影候选得分、保留候选集合和候选局部有界域\\
$\Delta\bm g^{\mathrm{ref}},d^{\mathrm{ref}}$ & 候选生成和局部投影采用的参考间隙增量与公共偏移\\
$\bm\theta^{\mathrm{cand}},\bm r^{\mathrm{cand}}$ & 投影搜索采用的完整模型候选点及其电压残差\\
$\widetilde{\bm r},\widetilde{\bm J}$ & 经 $\bm\Lambda_{\mathrm{VP}}^{1/2}$ 变换后的白化残差和灵敏度矩阵\\
$\mathcal S_g$ & 具有有效间隙标定信息、允许估计间隙增量的传感器集合\\
$\bm J_g,\bm J_{\nu}$ & 预测波形对间隙参数和空间/振动干扰参数的雅可比矩阵\\
$\eta_{g,s}$ & 第 $s$ 个间隙参数的局部条件可辨识性指标\\
$\delta g_s^{\mathrm{LS}}$ & 联合投影系统得到的第 $s$ 个局部线性间隙修正量\\
$\overline{\Delta g}_s$ & 投影修正后的局部间隙增量搜索中心\\
$h_{g,s}^{\mathrm{proj}}$ & 由条件投影残差与灵敏度之比得到的局部搜索尺度基准\\
$h_{g,s},h_{g,s}^{\mathrm{cons}}$ & 最终局部优化采用的间隙搜索半宽和预设保守半宽\\
\bottomrule
\end{longtable}

\clearpage
\small
\begin{thebibliography}{99}

\bibitem{Yu2020}
B. Yu, H. Ke, E. Shen, and T. Zhang,
``A review of blade tip clearance-measuring technologies for gas turbine engines,''
\emph{Measurement and Control}, vol. 53, no. 3--4, pp. 339--357, 2020.

\bibitem{Muller1997}
D. Muller, S. Mozumdar, E. Johann, and A. G. Sheard,
``Capacitive measurement of compressor and turbine blade tip to casing running clearance,''
\emph{Journal of Engineering for Gas Turbines and Power}, vol. 119, no. 4, pp. 877--884, 1997.

\bibitem{Golub1973}
G. H. Golub and V. Pereyra,
``The differentiation of pseudo-inverses and nonlinear least squares problems whose variables separate,''
\emph{SIAM Journal on Numerical Analysis}, vol. 10, no. 2, pp. 413--432, 1973.

\bibitem{Golub2003}
G. H. Golub and V. Pereyra,
``Separable nonlinear least squares: The variable projection method and its applications,''
\emph{Inverse Problems}, vol. 19, no. 2, pp. R1--R26, 2003.

\bibitem{Heller2020}
D. Heller, I. A. Sever, and C. W. Schwingshackl,
``A method for multi-harmonic vibration analysis of turbomachinery blades using Blade Tip-Timing and clearance sensor waveforms and optimization techniques,''
\emph{Mechanical Systems and Signal Processing}, vol. 142, Art. no. 106741, 2020.

\bibitem{Diamond2021}
D. H. Diamond, P. S. Heyns, and A. J. Oberholster,
``Constant speed tip deflection determination using the instantaneous phase of blade tip timing data,''
\emph{Mechanical Systems and Signal Processing}, vol. 150, Art. no. 107151, 2021.

\end{thebibliography}

\end{document}
